% !TEX root = ../discretemath.tex

\section{Многочлен Татта (корректность рекуррентного определения)}

\subsection{Определение}

\begin{Statement}{Удаление--стягивание для хроматического многочлена}{}
  Пусть $G$ --- граф, $e\in E(G)$ --- ребро, не являющееся петлёй. Тогда:
  \[
    \chi_G(x) = \chi_{G-e}(x) - \chi_{G \cdot e}(x),
  \]
  где $G-e$ --- граф с удалённым ребром $e$, а $G \cdot e$ --- граф,
  полученный стягиванием $e$ (отождествлением его концов).
\end{Statement}

\begin{Definition}{Мост}{}
  Ребро $e \in E(G)$ называется \emph{мостом}, если $c(G-e) > c(G)$,
  то есть его удаление увеличивает число компонент связности.
\end{Definition}

\begin{Proposition}{Удаление--стягивание для деревьев и лесов}{}
  Пусть $e$ --- не петля. Тогда:
  \[
    \#\{\text{ост.\ деревьев в }G\}
    = \#\{\text{ост.\ деревьев в }G-e\}
    + \#\{\text{ост.\ деревьев в }G \cdot e\},
  \]
  \[
    \#\{\text{ост.\ лесов в }G\}
    = \#\{\text{ост.\ лесов в }G-e\}
    + \#\{\text{ост.\ лесов в }G \cdot e\}.
  \]
\end{Proposition}

\begin{Definition}{Многочлен Татта}{}
  \emph{Многочлен Татта} $T_G(x,y)$ графа $G$ определяется рекурсией
  по выбранному ребру $e \in E(G)$:
  \begin{enumerate}
    \item[(i)]   если $e$ --- не петля и не мост:
                 $T_G(x,y) = T_{G-e}(x,y) + T_{G \cdot e}(x,y)$;
    \item[(ii)]  если $e$ --- петля:
                 $T_G(x,y) = y\, T_{G-e}(x,y)$;
    \item[(iii)] если $e$ --- мост:
                 $T_G(x,y) = x\, T_{G-e}(x,y)$;
    \item[(iv)]  если $G$ не имеет рёбер: $T_G(x,y) = 1$.
  \end{enumerate}
\end{Definition}

\begin{Example}{Вычисление $T(K_3)$}{}
  $K_3$ на вершинах $\{1,2,3\}$, рёбра $e_1=\{12\}$, $e_2=\{23\}$, $e_3=\{13\}$.
  Ребро $e_1$ --- не мост, не петля:
  \[
    T(K_3) = T(K_3 - e_1) + T(K_3 \cdot e_1).
  \]
  \begin{itemize}
    \item $K_3-e_1$: путь $1\text{-}3\text{-}2$, оба ребра --- мосты.
          $T = x \cdot x \cdot 1 = x^2$.
    \item $K_3 \cdot e_1$: вершины $\{12\},\{3\}$, два параллельных ребра.
          Одно из них не мост, не петля.
          $T = T(\text{мост}) + T(\text{петля}) = x + y$.
  \end{itemize}
  \[
    \boxed{T(K_3) = x^2 + x + y.}
  \]
\end{Example}

\begin{figure}[H]
	\centering
	\begin{tikzpicture}[
		>=Stealth, font=\small,
		v/.style={circle, draw=black, fill=black!75, minimum size=5pt, inner sep=0pt}]
		\begin{scope}[xshift=0cm]
			\node[v, label=above:{$1$}]       (A1) at ( 0,   0.9) {};
			\node[v, label=below left:{$2$}]  (A2) at (-0.8,-0.5) {};
			\node[v, label=below right:{$3$}] (A3) at ( 0.8,-0.5) {};
			\draw (A1)--(A2) (A2)--(A3);
			\draw[red, very thick] (A1) -- node[right, red, font=\scriptsize]{$e_1$} (A3);
			\node at (0,-1.4) {$K_3$};
		\end{scope}
		\draw[->, thick] (1.2,  0.3) -- (2.1,  0.9)
		node[midway, above left, font=\scriptsize] {$-e_1$};
		\draw[->, thick] (1.2, -0.2) -- (2.1, -0.9)
		node[midway, below left, font=\scriptsize] {$\cdot e_1$};
		\begin{scope}[xshift=3.5cm, yshift=1.5cm]
			\node[v, label=above:{$1$}] (B1) at (-0.7, 0.3) {};
			\node[v, label=above:{$3$}] (B2) at ( 0,  -0.4) {};
			\node[v, label=above:{$2$}] (B3) at ( 0.7, 0.3) {};
			\draw (B1)--(B2)--(B3);
			\node at (0,-1.0) {$K_3-e_1$};
			\node[blue!70!black] at (0,-1.5) {$T=x^2$};
		\end{scope}
		\begin{scope}[xshift=3.5cm, yshift=-2.0cm]
			\node[v, label=above:{$\{12\}$}] (C1) at (0, 0.5) {};
			\node[v, label=below:{$3$}]      (C2) at (0,-0.5) {};
			\draw[bend left=40]  (C1) to (C2);
			\draw[bend right=40] (C1) to (C2);
			\node at (0,-1.35) {$K_3\cdot e_1$};
			\node[blue!70!black] at (0,-1.9) {$T=x+y$};
		\end{scope}
		\node[blue!70!black, font=\bfseries] at (0,-2.4) {$T(K_3)=x^2+x+y$};
	\end{tikzpicture}
	\caption{Удаление--стягивание по $e_1$: путь даёт $x^2$ (два моста),
		параллельные рёбра дают $x+y$ (мост плюс петля после стягивания).}
\end{figure}

\subsection{Вспомогательные леммы}

\begin{Lemma}{Мультипликативность по компонентам}{}
  Если $G_1,\dots,G_k$ --- компоненты связности $G$, то
  $T_G = T_{G_1} \cdots T_{G_k}$.
\end{Lemma}

\begin{proof}
  Индукция по $|E(G)|$. База ($E=\varnothing$): $T_G=1$ и каждый $T_{G_i}=1$.

  Шаг: ребро $e\in E(G_1)$ имеет тот же тип (петля/мост/обычное) в $G$,
  что и в $G_1$. При любом правиле:
  \[
    G - e = (G_1-e)\sqcup G_2\sqcup\cdots,\quad
    G \cdot e = (G_1\cdot e)\sqcup G_2\sqcup\cdots.
  \]
  По ИП $T_{G-e} = T_{G_1-e}\prod_{j\ge 2}T_{G_j}$
  и $T_{G\cdot e} = T_{G_1\cdot e}\prod_{j\ge 2}T_{G_j}$.
  Выносим $\prod_{j\ge 2}T_{G_j}$ и получаем требуемое.
\end{proof}

\begin{Lemma}{Склейка по вершине}{}
  Пусть $G = G_1 \cup G_2$,\ $V(G_1)\cap V(G_2)=\{v\}$,
  $E(G_1)\cap E(G_2)=\varnothing$. Тогда $T_G = T_{G_1}\,T_{G_2}$.
\end{Lemma}

\begin{proof}
  Индукция по $|E(G_1)|+|E(G_2)|$. Если $E(G_1)=\varnothing$, то $G=G_2$.

  Возьмём $e\in E(G_1)$. Тип $e$ в $G$ совпадает с типом в $G_1$:
  приклейка $G_2$ в вершину $v$ не создаёт обходных путей внутри $G_1$.
  При любом правиле $G-e=(G_1-e)\cup G_2$ и $G\cdot e=(G_1\cdot e)\cup G_2$.
  По ИП $T_{G\pm e}=T_{G_1\pm e}\,T_{G_2}$, откуда при подстановке в рекурсию:
  \[
    T_G = (T_{G_1-e}+T_{G_1\cdot e})\,T_{G_2} = T_{G_1}\,T_{G_2}.
  \]
\end{proof}

\subsection{Корректность определения}

\begin{Proposition}{Корректность}{}
  $T_G(x,y)$ не зависит от порядка применения правил удаления--стягивания.
\end{Proposition}

\begin{proof}
  Достаточно показать, что результат не меняется при перестановке
  двух соседних рёбер $e_i$ и $e_{i+1}$ в некотором фиксированном порядке
  (любая перестановка --- произведение соседних транспозиций).

  Пусть $G'$ --- граф после обработки $e_1,\dots,e_{i-1}$.
  Рассматриваем все возможные случаи для пары $(e_i,e_{i+1})$ в $G'$.

  \subsubsection*{Случаи 0 и 1: кратные рёбра или петля/мост}

  Если $e_i,e_{i+1}$ кратны --- одинаковые графы в обоих порядках.
  Если одно из них --- петля или мост --- множитель $y$ или $x$ коммутирует
  со всеми операциями над другим ребром. Во всех этих случаях порядок не важен.

  \subsubsection*{Случай 2: оба не мосты, не петли}

  \noindent\textbf{Подслучай 3.1.} $c(G'-e_i-e_{i+1})=c(G')$.

  Каждое из рёбер остаётся не мостом после удаления другого.
  Раскрывая рекурсию в обоих порядках и используя коммутативность
  операций удаления и стягивания для различных рёбер:
  \[
    (G'\cdot e_j)-e_k \cong (G'-e_k)\cdot e_j,\quad
    (G'\cdot e_j)\cdot e_k \cong (G'\cdot e_k)\cdot e_j,
  \]
  получаем одну и ту же сумму из четырёх слагаемых.

  \medskip
  \noindent\textbf{Подслучай 3.2.} $c(G'-e_i-e_{i+1})=c(G')+1$.

  Оба ребра вместе образуют разрез; по отдельности --- нет.
  По лемме о мультипликативности сводим к компоненте $H$,
  содержащей оба ребра. Граф $G'-e_i-e_{i+1}$ имеет ровно
  две компоненты $A$ и $B$; ребро $e_i$ ведёт из $A$ в $B$,
  аналогично $e_{i+1}$.

  \begin{figure}[H]
  \centering
  \begin{tikzpicture}[>=Stealth, font=\small]
    \draw[fill=blue!10, draw=blue!50, thick, rounded corners=8pt]
      (-2.4,-1.1) rectangle (-0.3, 1.1);
    \node at (-1.35, 0) {$A$};
    \draw[fill=red!10, draw=red!50, thick, rounded corners=8pt]
      ( 0.3,-1.1) rectangle ( 2.4, 1.1);
    \node at ( 1.35, 0) {$B$};
    \node[circle, fill=blue!60, minimum size=5pt, inner sep=0pt] (a1) at (-0.3, 0.4){};
    \node[circle, fill=red!60,  minimum size=5pt, inner sep=0pt] (b1) at ( 0.3, 0.4){};
    \draw[blue!70!black, very thick] (a1)--node[above,font=\scriptsize]{$e_i$}(b1);
    \node[circle, fill=blue!60, minimum size=5pt, inner sep=0pt] (a2) at (-0.3,-0.4){};
    \node[circle, fill=red!60,  minimum size=5pt, inner sep=0pt] (b2) at ( 0.3,-0.4){};
    \draw[red!70!black,  very thick] (a2)--node[below,font=\scriptsize]{$e_{i+1}$}(b2);
  \end{tikzpicture}
  \caption{Случай 3.2: компоненты $A$ и $B$ связаны только рёбрами $e_i$
    и $e_{i+1}$. По отдельности каждое --- не мост; оба вместе --- разрез.}
  \end{figure}

  После стягивания $e_{i+1}$ граф $(G'\cdot e_{i+1})-e_i$ является
  $1$-суммой $A$ и $B$, поэтому по лемме о склейке
  $T_{(G'\cdot e_{i+1})-e_i}=T_A\cdot T_B$.
  Симметрично $T_{(G'\cdot e_i)-e_{i+1}}=T_A\cdot T_B$.
  Стягивания коммутируют:
  $T_{(G'\cdot e_{i+1})\cdot e_i}=T_{(G'\cdot e_i)\cdot e_{i+1}}$.

  Итого раскрытие сначала по $e_{i+1}$:
  \[
    T_{G'}=x\,T_A T_B + T_A T_B + T_{(G'\cdot e_i)\cdot e_{i+1}},
  \]
  и по $e_i$ даёт то же выражение (симметрично) --- результат совпадает.
\end{proof}
